\begin{question}

  \begin{questionpart}
    Find the linear combination of $\ket{+}$ and $\ket{-}$ kets that
    maximizes the uncertainty product
    \begin{equation*}
      \dispersion{S_x}\dispersion{S_y}
    \end{equation*}
  \end{questionpart}
  
  \begin{questionpart}
    Verify explicitly that for the linear combination you found, the
    uncertainty relation for $S_x$ and $S_y$ is not violated.
  \end{questionpart}
  
\end{question}

\begin{purpose}
  Pretty sure Sakurai wants us to connect our previous problem in
  which we derived the generic ket-on-a-swivel expression so we can
  show the angle that maximizes the uncertainty outcome.  We may not
  need the eigenvectors that we derived in problem
  \ref{1.11:theQuestion}, but we can at least rob some of the math.
\end{purpose}

\begin{answer}

  \begin{answerpart}{Maximus Dipersionus Uncertainus}

    \begin{equation}
      \begin{split}
        \dispersion{S_x}\dispersion{S_y}
        &=
        \left(\expval{S_x^2} - \expval{S_x}^2\right)
        \left(\expval{S_y^2} - \expval{S_y}^2\right)
        \\
      \end{split}
    \end{equation}

    We can simplify our characteristic equation above by recalling
    (from any number of previous problems) that for all the spin
    operators $\expval{S_n^2} = \frac{\hbar^2}{4}$
    
    \begin{equation}
      \inlabel{top}
      \begin{split}
        \dispersion{S_x}\dispersion{S_y}
        &=
        \left(\frac{\hbar^2}{4} - \expval{S_x}^2\right)
        \left(\frac{\hbar^2}{4} - \expval{S_y}^2\right)
        \\
      \end{split}
    \end{equation}

    There's a subtlety here which is worth calling out.  The problem
    isn't looking for the linear combination of $\ket{+}$ and
    $\ket{-}$ with \emph{real} coefficients.  It's asking for the
    combination from the entire Bloch Sphere.\footnote{We haven't
    covered the Bloch Sphere yet but it's worth looking up quickly.}
    Blessedly, we've already computed the eigenvectors for a linear
    combination of operators.  In problem \ref{1.11:theQuestion} we
    computed the eigenvectors of a spin operator which is rotated
    arbitrarily in 3D space by two angles $\alpha$ (azimuthal angle in
    the $xy$ plane starting from the $x$ axis) and $\beta$ (elevation
    angle starting from the $z$ axis).  Restating equation
    \ref{1.11:spinAngle}, we found the eigenvectors of this
    arbitrarily-rotated spin operator to be:

    \begin{equation}
      \inlabel{spinAngle}
      \ket{\alpha_+} =
      \cos{\frac{\beta}{2}} \ket + + e^{i\alpha}\sin{\frac{\beta}{2}} \ket -
    \end{equation}

    OK...what does that mean exactly?  The expression above is the
    generalized expression of the $\ket{S_n+}$ eigenket from the
    generalized operator $S_n$ which has been arbitrarily rotated in
    the Bloch Sphere.  So the question then becomes...how does that
    help us?  Let's start by going back to the setup for problem
    \ref{1.11:theQuestion}

    \begin{equation}
      \vectorbold{S}\cdot\vectorbold{\hat{n}} \ket{\vectorbold{S}\cdot\vectorbold{\hat{n}}; +}
      = \left(\frac{\hbar}{2}\right)\ket{\vectorbold{S}\cdot\vectorbold{\hat{n}}; +}
    \end{equation}

    So...does that help us?  No.  No it does not.  So, like we said
    earlier, the eigen-nature of our previous finding doesn't really
    help us, but we \emph{can} use the resulting eigen-ket because it
    has all the properties we care about.  So if we move to the next
    dead-end (well not so much dead-end as bad-idea).  We can use our
    handy expression from equation \inref{spinAngle} in our earlier
    equation \inref{top}.  Then we can maximize the value by checking
    partial derivatives.

    \begin{equation}
      \inlabel{top}
      \begin{split}
        0
        &= \pdv{}{\alpha}{\beta} \dispersion{S_x}\dispersion{S_y}\\
        &= \pdv{}{\alpha}{\beta}
        \left(\frac{\hbar^2}{4} - \expval{S_x}^2\right)
        \left(\frac{\hbar^2}{4} - \expval{S_y}^2\right)
        \\
        &= \pdv{}{\alpha}{\beta}
        \left(\frac{\hbar^2}{4} - \mel{\alpha+}{S_x}{\alpha+}^2\right)
        \left(\frac{\hbar^2}{4} - \mel{\alpha+}{S_y}{\alpha+}^2\right)
        \\
        &= \pdv{}{\alpha}{\beta}
        \left(\frac{\hbar^2}{4}
        [- (\cos{\frac{\beta}{2}} \ket + + e^{i\alpha}\sin{\frac{\beta}{2}} \ket -)
        S_x
        (\cos{\frac{\beta}{2}} \ket + + e^{i\alpha}\sin{\frac{\beta}{2}} \ket -)]^2\right)
        \left(\frac{\hbar^2}{4}
        \ldots
        \right)
        \\
      \end{split}
    \end{equation}

    So...yeah...be my guest...I don't even know if I did those
    substitutions properly, and I don't care.  Let's try this another
    way then, shall we?

    Going back to \inref{top}, we see that we can rewrite it as

    \begin{equation}
      \begin{split}
        \dispersion{S_x}\dispersion{S_y} = 
        &\left[(\text{positive real number}) - (\text{another positive real number})\right]\\
        &\left[(\text{positive real number}) - (\text{another positive real number})\right]\\
      \end{split}
    \end{equation}

    So, if we're trying to maximize that expression, then we probably
    want the two numbers being multiplied are maximized.  We can do
    that by ensuring that the second number there is minimized.  Let's
    see what we can do to make that happen.  Ideally, we'll be able to
    make them both 0.  So let's see if we can make that happen, we're
    going to be looking for $\mel{\psi}{S_x}{\psi} = 0$ and
    $\mel{\psi}{S_y}{\psi} = 0$.

    \begin{equation}
      \tag{1.111a}
      S_x = \frac{\hbar}{2} [\ket - \bra + + \ket + \bra - ]
    \end{equation}
    
    \begin{equation}
      \begin{split}
        0
        &= \mel{\psi}{S_x}{\psi}
        \\
        &=
        \left\{\cos{\frac{\beta}{2}} \bra + + e^{-i\alpha}\sin{\frac{\beta}{2}} \bra -\right\}
        \left\{\frac{\hbar}{2} [\ket - \bra + + \ket + \bra - ]\right\}
        \left\{\cos{\frac{\beta}{2}} \ket + + e^{i\alpha}\sin{\frac{\beta}{2}} \ket -\right\}
        \\
        &=
        \frac{\hbar}{2}
        \left\{\cos{\frac{\beta}{2}} \bra + + e^{-i\alpha}\sin{\frac{\beta}{2}} \bra -\right\}
        \left\{\cos{\frac{\beta}{2}} \ket - + e^{i\alpha}\sin{\frac{\beta}{2}} \ket +\right\}
        \\
        &=
        \frac{\hbar}{2}\left\{
        e^{i\alpha}\sin{\frac{\beta}{2}}\cos{\frac{\beta}{2}}\braket{+}{+}
        +
        e^{-i\alpha}\sin{\frac{\beta}{2}}\cos{\frac{\beta}{2}}\braket{-}{-}
        \right\}
        \\
        &=
        \frac{\hbar}{2}
        \sin{\frac{\beta}{2}}\cos{\frac{\beta}{2}}
        \left(e^{i\alpha} + e^{-i\alpha}\right)
        \\
        &=
        2\frac{\hbar}{2}
        \sin{\frac{\beta}{2}}\cos{\frac{\beta}{2}}
        \cos{\alpha}
        \\
        &=
        \frac{\hbar}{2}
        \sin{\beta}
        \cos{\alpha}
        \\
      \end{split}
    \end{equation}

    So, we know from $S_x$ that so long as $\beta$ is an integer
    multiple of $\pi$ or $\pi/2$ for $\alpha$ then $S_x$ will be 0.
    Let's take a look at $S_y$

    \begin{equation}
      \tag{1.111b}
      S_y = i \frac{\hbar}{2} [\ket - \bra + - \ket + \bra - ]
    \end{equation}

    \begin{equation}
      \begin{split}
        0
        &= \mel{\psi}{S_y}{\psi}
        \\
        &=
        \left\{\cos{\frac{\beta}{2}} \bra + + e^{-i\alpha}\sin{\frac{\beta}{2}} \bra -\right\}
        \left\{i \frac{\hbar}{2} [\ket - \bra + - \ket + \bra - ]\right\}
        \left\{\cos{\frac{\beta}{2}} \ket + + e^{i\alpha}\sin{\frac{\beta}{2}} \ket -\right\}
        \\
        &=
        i\frac{\hbar}{2}
        \left\{\cos{\frac{\beta}{2}} \bra + + e^{-i\alpha}\sin{\frac{\beta}{2}} \bra -\right\}
        \left\{\cos{\frac{\beta}{2}} \ket - - e^{i\alpha}\sin{\frac{\beta}{2}} \ket +\right\}
        \\
        &=
        i\frac{\hbar}{2}\left\{
        e^{-i\alpha}\sin{\frac{\beta}{2}}\cos{\frac{\beta}{2}}\braket{-}{-}
        -
        e^{i\alpha}\sin{\frac{\beta}{2}}\cos{\frac{\beta}{2}}\braket{+}{+}
        \right\}
        \\
        &=
        -i\frac{\hbar}{2}
        \sin{\frac{\beta}{2}}\cos{\frac{\beta}{2}}
        \left(e^{i\alpha} - e^{-i\alpha}\right)
        \\
        &=
        2\frac{\hbar}{2}
        \sin{\frac{\beta}{2}}\cos{\frac{\beta}{2}}
        \sin{\alpha}
        \\
        &=
        \frac{\hbar}{2}
        \sin{\beta}
        \sin{\alpha}
        \\
      \end{split}
    \end{equation}

    Clearly, any solution where $\beta = n\pi$ will result in the
    maximum uncertainty.  If we reference our figure
    \ref{1.11:spinAxes} (or figure \ref{q:spinAxes} in the \ac{qrf}).
    You'll note also that the angle $\alpha$ doesn't really help us,
    and that the assigned $\beta$ leaves our vector colinear with the
    $z$ axis.  That brings us to the \emph{easy} way to solve this
    problem (have you noticed how I always leave the easy way for the
    very end?).

    If you want the \emph{maximum uncertainty} then you want the
    \emph{minimum knowledge}.  When you prepare a state that is in the
    $z$ axis (either plus or minus) then you have \emph{absolutely no
    knowledge} about its orientation in the $xy$ plane.  When you have
    no knowledge, then you're just as likely to get a $\ket{x/y+}$ as
    you are to get an $\ket{x/y-}$ so the expectation value is always
    going to be 0.  That's the easy way to solve this part.  But take
    heart, we're going to need the math above for the next part...or
    are we?  No, we aren't...but aren't you glad you did it anyway?

    Oh, and because we're technically required to give the linear
    combination, I present unto you:

    \begin{equation}
      \ket{\alpha} = \ket{\pm}\qed
    \end{equation}
    
  \end{answerpart}
  
  \begin{answerpart}{But are you \emph{suuuuure}?}
    So, now that we want to explicitly verify the relation, we can
    just plug our 0's into equation \inref{top} and find that the
    product of the dispersions is $\hbar^4/16$.  So let's see how that
    stacks up against the commutation.

    \begin{equation}
      \begin{split}
        [S_x, S_y]
        &=  S_xS_y - S_yS_x\\
        &= i\frac{\hbar^2}{4}
        \left(
             [\ket - \bra + + \ket + \bra - ]
             [\ket - \bra + - \ket + \bra - ]
             - 
             [\ket - \bra + - \ket + \bra - ]
             [\ket - \bra + + \ket + \bra - ]
             \right)
        \\
        &= i\frac{\hbar^2}{4}
        \left(
             [\dyad{+}{+} -\dyad{-}{-}]
             - 
             [\dyad{-}{-} - \dyad{+}{+}]
             \right)
        \\
        &= 2i\frac{\hbar^2}{4}(\Lambda_+ - \Lambda_-)\\
        &= 2i\frac{\hbar}{2}S_z\\
      \end{split}
    \end{equation}

    So, if we evaluate our final equation

    \begin{equation}
      \begin{split}
        \frac{1}{4}\abs{\expval{[S_x, S_y]}}^2
        &=
        \frac{1}{4}
        \abs{
          \mel{\pm}{2i\frac{\hbar}{2}S_z}{\pm}
        }^2
        \\
        &=
        \abs{
          i\frac{\hbar^2}{4}\braket{\pm}{\pm}
        }^2
        \\
        &= \frac{\hbar^4}{16} \\
      \end{split}
    \end{equation}

    \begin{equation}
      \begin{alignedat}{2}
        & &
        \dispersion{S_x}\dispersion{S_y}
        &\overset{?}{\ge}
        \frac{1}{4}\abs{\expval{[S_x, S_y]}}^2\\
        \implies \quad & &
        \frac{\hbar^4}{16} &\ge \frac{\hbar^4}{16}\qed\\
      \end{alignedat}
    \end{equation}
    
  \end{answerpart}

\end{answer}
